\subsection{Identifying Irreducible Quadratics}
We already have a tool to determine whether a quadratic is factorable (with rational factors), prime, or irreducible.
\begin{definition}{Discriminant}
The \term{discriminant} of a quadratic equation $ax^2+bx+c=0$ is the number
$$D=b^2-4ac$$
whose square root appears in the quadratic formula. The quadratic equation has:
\begin{itemize*}
\item Two real solutions if $D>0$
\item One real solution if $D=0$
\item No real solutions if $D<0$
\end{itemize*}
If $D$ is a perfect square, the solutions are rational. Otherwise, they are irrational.
\end{definition}
Putting this in terms of factorability:
\begin{itemize}
\item If $D$ is a \makeblank, the quadratic has factors with \makeblank[1.5in] coefficients.
\item Otherwise, if $D$ is \makeblank[1.5in], the quadratic is prime (but has factors with \makeblank[1.5in] coefficients).
\item If $D$ is \makeblank[1.5in], the quadratic is irreducible.
\end{itemize}
\begin{Example}
Identify each quadratic expression as factorable (with rational factors), prime, or irreducible.
\begin{enumerate}\setabc
\item $x^2-6x+4$\makeblanknewf
\item $x^2-6x+8$\makeblanknewf
\item $x^2-6x+13$\makeblanknewf
\end{enumerate}
\end{Example}